\exnumber{00}
\extitle{Projet 3 : Let's fair trade}
\exfiles{ex00/archive zip, nomme Name\_FirstName\_DeliverableName\_Date\_VersionNumber}
\exauthorize{Ceux nommés dans l'exercice}
\makeheaderfiles

Création d'une plateforme e-commerce de vêtements d'occasion sur desktop et mobile

Bruno adore la mode. Il adore acheter des vêtements ! Comme il est soucieux de l'environnement, il préfère acheter des vêtements d'occasion. Cependant, les plateformes existantes ne sont pas très user-friendly.

« J'ai besoin d'en savoir plus sur l'impact environnemental des pièces que j'achète. J'aimerais savoir quel âge elles ont, d'où elles viennent, pour mesurer leur durée de vie et le nombre de kilomètres qu'elles ont parcourus. Quand un nouveau jean arrive dans le magasin de prêt-à-porter en Europe, il a consommé environ 11 000 litres d'eau et parcouru 65 000 kilomètres. J'aimerais suivre ce type de données. »

Vous réunissez une équipe de 3 pour réfléchir et concevoir un site web et une application pour des produits d'occasion.

Vous devez rencontrer Bruno après 2 semaines pour une première présentation de votre travail, après avoir réalisé un prototype mid-fidelity testé.

La présentation finale de votre travail aura lieu 2 semaines après la présentation intermédiaire, où vous montrerez un prototype haute fidélité testé.

\textbf{Durée estimée :} 4 semaines

\textbf{Taille de l'équipe :} 3 personnes max

\textbf{Consignes :}

\begin{itemize}
\item Posez une stratégie UX et UI avant de foncer. Remettez le brief initial en question avec une recherche utilisateur.
\item Respectez les étapes : empathie, définir, énoncé du problème, idéation, prototype, test.
\item Vous devez interroger de vraies personnes pendant l'empathie, la recherche et les tests.
\item Dans les deux phases, testez vos prototypes avec au moins 5 personnes.
\item Faire une présentation comme si vous montriez votre travail à Bruno. Assurez-vous qu'il y a au moins 5 personnes dans l'audience.
\item Déposez les documents sur Github, dans un zip nommé « ProjectTitle\_Name\_FirstName », et chaque livrable ainsi : « Name\_FirstName\_DeliverableName\_Date\_VersionNumber ». Exemple : « Doe\_John\_MidFiPrototype\_05242024\_V1 ».
\item Ajoutez un titre dans tous les documents écrits.
\item Soignez l'aspect global des livrables. Ils doivent être clairs, simples et lisibles. Vous pouvez regarder canva.com, sans surcharger le dessin de détails.
\item Préparez un formulaire de retour pour le public. Il contient au moins les questions listées plus bas.
\item Dans les deux présentations, vous parlez à Barbara et à son équipe.
\item Durée : 10 minutes de présentation et 5 minutes de questions.
\item Vous pouvez vous servir des audits de présentation pour structurer vos présentations.
\item Amusez-vous !
\end{itemize}

\textbf{Outils :}

\begin{itemize}
\item Google Forms \url{https://docs.google.com/forms/u/0/}
\item Airtable \url{https://www.airtable.com/}
\item Typeform \url{https://www.typeform.com/}
\end{itemize}

\textbf{Pistes :}

\begin{itemize}
\item Compétences de présentation efficaces \url{https://academic.oup.com/femsle/article/364/24/fnx235/4587905}
\end{itemize}

Livrables intermédiaires :

\textbf{Citation :}

« Your mission in any presentation is to inform, educate, and inspire. » Robert Ballard

\medskip

\textbf{Phase 1 : audit de la présentation intermédiaire}

Vérifiez :

\begin{itemize}
\item Au moins 5 personnes assistent à la présentation
\item La présentation porte sur tout le processus UX, y compris :
\begin{itemize}
\item Le brief initial
\item Un planning, avec les phases UX, les outils choisis, le calendrier
\item Les résultats de la recherche utilisateur, qualitative et quantitative (persona, chiffres, etc.)
\item Un énoncé de problème qui ajuste le brief initial d'après la recherche
\item Une introduction du concept sorti de l'idéation
\item Le prototype moyenne fidélité animé (une vidéo de démo)
\item Les retours principaux du test d'utilisabilité
\item Les prochaines étapes du projet
\end{itemize}
\item La présentation ne dépasse pas 10 minutes
\item Les formulaires de retour ont été distribués au début de la présentation
\item Vous vous êtes amusés !
\end{itemize}

\medskip

\textbf{Phase 2 : audit de la présentation finale}

Vérifiez :

\begin{itemize}
\item Au moins 5 personnes assistent à la présentation
\item La présentation porte sur le processus UI, y compris :
\begin{itemize}
\item Un bref rappel de la présentation précédente
\item Un planning, avec les phases UI, les outils choisis, le calendrier
\item Une planche d'ambiance
\item Un styleguide (couleurs, marque, typographie, quelques éléments de la bibliothèque)
\item Le prototype haute fidélité animé (une vidéo de démo)
\item Les retours principaux des tests d'utilisabilité et de désirabilité
\item Les prochaines étapes du projet
\end{itemize}
\item La présentation ne dépasse pas 10 minutes
\item Les formulaires de retour ont été distribués au début de la présentation
\item Vous vous êtes amusés !
\end{itemize}

\textbf{Questions à poser dans les formulaires de retour}

\begin{itemize}
\item Que retenez-vous de cette présentation ?
\item De 0 à 10, le discours était-il clair ?
\item Pourquoi ?
\item De 0 à 10, la présentation était-elle engageante ?
\item Pourquoi ?
\item De 0 à 10, la présentation était-elle cohérente ?
\item Pourquoi ?
\item De 0 à 10, les visuels portaient-ils ?
\item Pourquoi ?
\item De 0 à 10, le problème était-il clair ?
\item Pourquoi ?
\item De 0 à 10, le prototype répond-il au problème ?
\item Pourquoi ?
\item Quelle recommandation feriez-vous pour améliorer cette présentation ?
\item Pourquoi ?
\end{itemize}

